\section{Involutions of the three-cusp group}

Throughout this part \(F=\widehat F(x,y)\), \(z=(xy)^{-1}\), and
\(\mathcal P\) denotes the subgroup of \(\operatorname{Out}(F)\)
preserving each named peripheral conjugacy class up to a common
unit power. All group homomorphisms are continuous. We use the
normalized lift
\[
 (\lambda,f):\quad x\longmapsto x^\lambda,\qquad
 y\longmapsto f^{-1}y^\lambda f,\qquad f\in F',
\]
where \(F'\) is the closed derived subgroup. This normalization is
a homomorphic section, as in the refereed LM4 D.32. Write
\(c=(-1,1)\). We use R1--R3 only with their stated scopes.

\statement{A.1}{Known}{Torsion in an acting group and its action on torsion are different assertions}
Lochak and Schneps (1997, Proposition 4(i)) prove
\(C_F(c)=1\); Proposition 4(ii) proves that \(c\) is
self-centralizing in \(\widehat{GT}\).
Schneps (1997, \S1.4(4)) and Lochak and Schneps (2006,
Question 1.4) ask whether all nonidentity torsion in
\(\widehat{GT}\) consists of conjugates of \(c\).
Harbater and Schneps (2000, Main Theorem(a), proved by
Propositions 3--4 and Theorem 2) identify
\(\widehat{GT}_0\) with the centralizer in \(\mathcal P\) of
the geometric \(S_3\).
Collas's genus-zero Theorems 3.3 and 3.10 identify prime-order
torsion in \(\widehat\Gamma_{0,[n]}\) with discrete torsion
and prove the \(\lambda\)-conjugacy action of \(\widehat{GT}\).
His genus-one Theorems 6.3 and 6.5 give the corresponding
statements with acting group \(\widehat{GR}\), isomorphic
to the subgroup denoted \(I\Gamma\) in that article;
replacing it there by all of \(\widehat{GT}\) is not
justified by those theorems.\cite{A_LS1997,A_Schneps1997,A_LS2006,A_HS2000,A_Collas2012a,A_Collas2012b}

\begin{proof}
For the centralizer comparison identify \(K(0,4)\) with
\(\widehat F(x,y)\). The natural \(S_4\)-action on its outer
automorphism group factors through the anharmonic \(S_3\).
Harbater--Schneps's quasi-special condition is precisely the
three peripheral power condition here; their symmetric
condition is commutation with that \(S_3\). Their proof
translates commutation with the transposition into the
two-cycle relation and commutation with the three-cycle into
the hexagon. The identification has the required source
and target, not merely an abstract isomorphism.

In Collas's formulas the variable of prime order lies in a
completed mapping class group; \((\lambda,f)\) is an arbitrary
element of the acting group. The conclusion
\(F(\gamma)=u\gamma^\lambda u^{-1}\) does not classify
finite-order elements \(F\) of the acting group. The subgroup
\(\widehat{GR}\) in the genus-one article is retained in its
Theorem 6.5. Neither action theorem therefore proves task A.8.
\end{proof}

The opened sources establish no literature antecedent for the
full order bound in R1. Searches included the authors' source
lists, the cited torsion papers, and recent work on profinite
curve-complex automorphisms and Galois--Teichm\"uller theory.
This records the search, not a priority claim or a proof that
no antecedent exists.

\statement{A.2}{Proved}{Relative pronilpotent levels recover the full peripheral group}
For an open characteristic subgroup \(U\) of \(F\), put
\[
 R_U=R_{\mathrm{nil}}(U),\quad F_U=F/R_U,\quad
 P_U=\operatorname{Out}^{\sharp}(F_U).
\]
For \(V\subset U\), the quotient \(F_V\to F_U\) induces
a continuous epimorphism \(P_V\to P_U\). There is an
isomorphism of profinite groups
\[
                 \mathcal P \xrightarrow{\sim}\varprojlim_U P_U.
\]
Two elements of \(\mathcal P\) are conjugate in \(\mathcal P\)
if and only if their images are conjugate in every \(P_U\).
This proves Lead 1.

\begin{proof}
Restricting finite nilpotent quotients of \(U\) to \(V\)
shows \(R_V\subset R_U\). Each \(R_U\) is characteristic in
\(F\). Since \(R_U\subset U\) and open characteristic
subgroups are cofinal in the finitely generated profinite
group \(F\), we have \(\bigcap_U R_U=1\).

We first prove surjectivity onto \(P_U\). Choose a representative
of a peripheral outer class with common multiplier \(\lambda\),
and a Galois element with cyclotomic multiplier \(\lambda\);
the cyclotomic character is onto
\(\widehat{\mathbb Z}^{\times}\). The latter acts on \(F_U\).
Composing with its inverse gives a multiplier-one peripheral
automorphism. R2 lifts this automorphism, with chosen
conjugators, to \(F\). Composing back proves surjectivity.
In particular every peripheral automorphism of \(F_V\)
lifts to \(F\), so preserves the image of \(R_U\).
This proves that the transition maps are defined on the
entire displayed outer groups. Their surjectivity follows
by lifting first to \(\mathcal P\).

Here the automorphism groups are compact: intersecting
kernels of all maps to finite groups of bounded order
gives a cofinal system of open fully invariant subgroups.
The conditions that a continuous endomorphism induces
bijections on these finite quotients are closed conditions
on its generator images. Consequently \(\operatorname{Aut}(F)\)
is profinite. Inner automorphisms form a compact subgroup.
The peripheral conditions are closed, since conjugators
and the common multiplier range over compact groups.
The same argument applies to the finitely generated
quotients \(F_U\).

If \([a]\in\mathcal P\) dies in every \(P_U\), the sets
\[
 \{g\in F:a(x)\equiv gxg^{-1},\
              a(y)\equiv gyg^{-1}\pmod {R_U}\}
\]
are nonempty compact sets with the finite intersection
property. Their intersection gives an actual inner
representative for \(a\), proving injectivity.
For a compatible family of classes in \(P_U\), their
nonempty compact inverse images in \(\mathcal P\) have
the same finite intersection property, proving surjectivity.
Finally apply this argument to the transporter sets
\(\{g:g_Ua_Ug_U^{-1}=b_U\}\) in \(\mathcal P\).
Surjectivity onto \(P_U\) makes these nonempty whenever
the two images are conjugate. Their intersection is the
required full conjugator.
\end{proof}

\statement{A.3}{Proved}{A central quotient is the exact obstruction to a peripheral lift}
Let \(N\) be a closed characteristic subgroup of \(F\), and
let \([N,F]\) denote the closed subgroup generated by
commutators between \(N\) and \(F\).
An automorphism \(a\) of \(F/N\), with specified conjugators
\(b_i\) satisfying
\[
 a(\bar x_i)=b_i\bar x_i b_i^{-1},\qquad
 (x_1,x_2,x_3)=(x,y,z),
\]
lifts to a peripheral automorphism of \(F\) with these
conjugators if and only if it lifts to a peripheral
\emph{automorphism} of \(F/[N,F]\) with conjugators above
the \(b_i\). This proves Lead 2.

\begin{proof}
Necessity follows from characteristicity. For sufficiency
let \(a^\#\) and \(b_i^\#\) be the data on \(F/[N,F]\).
We use the finite central-refinement count in the refereed
LM7\_PUSH D.50, equations (D.50.6)--(D.50.16).
For any actual quotient \(q_E:F\twoheadrightarrow E\)
and normal subgroup \(A\triangleleft E\), that count
constructs an actual quotient
\(q^\dagger:F\twoheadrightarrow D^\dagger\), with central
kernel over \(D=E/A\). A second epimorphism to
\(D^\dagger\) with specified peripheral conjugators has
the same number of epimorphism lifts to \(E\), with the
transported \(A\)-conjugacy classes on the boundary.
This counting assertion uses only normality of \(A\).
The hypothesis \([E,A]=A\) occurs in the separate
joint-image assertion (D.50.17), not used here.

Take \(A=q_E(N)\). The image \(q^\dagger(N)\) is central,
so \(q^\dagger\) kills \([N,F]\).
Thus \(q^\dagger a^\#\) is an epimorphism from the
same source to \(D^\dagger\), with conjugators the
images of \(b_i^\#\). The finite count gives an epimorphism
to \(E\) whose conjugators agree with the \(b_i\) in \(D\);
the reference count is positive since it includes \(q_E\).

In the compact space \(F^3\) of conjugators impose the
boundary product relation, generation in every finite
quotient, and agreement with each \(b_i\) in every finite
quotient of \(F/N\). A finite list of conditions lies in
one joint finite quotient \(E\), and the preceding
construction satisfies that list. Compactness gives
full conjugators satisfying all conditions. Their
endomorphism is surjective since its compact image
maps onto every finite quotient; Hopficity makes it
an automorphism. Its generator images induce exactly
\(a\) modulo \(N\).
\end{proof}

\statement{A.4}{Proved}{A degree-three coefficient imposed by the third cusp}
Over a field \(k\) of characteristic zero or greater
than three, use the free class-three Lie algebra on
\(X,Y\), and set
\[
 Z=[X,Y],\quad U_3=[X,Z],\quad V_3=[Y,Z],\quad
 x=\exp X,\quad y=\exp Y .
\]
If \(x\mapsto x^\mu,\ y\mapsto g^{-1}y^\mu g\), with
\(\mu\in k^\times\) and
\(\log g=tZ+\text{central terms}\), sends \(xy\)
to a conjugate of \((xy)^\mu\), then
\(t=(\mu^2-1)/24\).

\begin{proof}
The Baker--Campbell--Hausdorff formula gives
\[
 W=\log(xy)=X+Y+\tfrac12Z+\tfrac1{12}(U_3-V_3).
\]
The logarithm of the image is
\[
 \mu(X+Y)+\tfrac{\mu^2}{2}Z+
       \tfrac{\mu^3}{12}(U_3-V_3)+\mu tV_3 .
\]
Write the logarithm of a conjugator as
\(aX+bY+dZ+\text{central terms}\).
In \(\exp(\operatorname{ad}(aX+bY+dZ))(\mu W)\),
comparison of degree two gives \(m=a-b=(\mu-1)/2\).
The difference of the \(U_3\)- and \(V_3\)-coefficients,
divided by \(\mu\), then gives
\[
 \frac{\mu^2}{6}-t=\frac16+\frac m2+\frac{m^2}{2}
                  =\frac{3\mu^2+1}{24}.
\]
The result follows. The coefficient \(d\) and all
central coordinates cancel.
\end{proof}

\statement{A.5}{Refuted}{The specified conjugator in task A.3 need not preserve the third cusp}
There is an involution \(\alpha\in\widehat{GT}\), conjugate
to \(c\) by a Galois element, for which the unique solution
\(h\in F'\) of \(f_\alpha=h^{-1}\alpha(h)\) gives
\[
                 \psi(x)=x,\qquad\psi(y)=hyh^{-1}
\]
that does not preserve the conjugacy class of \(z\)
on \(F/R_{\mathrm{nil}}(F)\). The failure occurs in
the class-three exponent-seven quotient of order \(7^5\).

\begin{proof}
Choose \(\theta=(\mu,g)\) in the Galois image with
\(\mu\equiv2\pmod7\), possible by surjectivity of the
cyclotomic character modulo seven, and put
\(\alpha=\theta c\theta^{-1}\).
Normalization is a homomorphic section, so these
are actual automorphisms with that equality.
From \(\alpha\theta(y)=\theta(y)^{-1}\) we obtain
\[
 \alpha(y)=\alpha(g)g^{-1}y^{-1}g\alpha(g)^{-1},
 \qquad f_\alpha=g\alpha(g)^{-1}.
\]
Thus \(h=g^{-1}\) is a solution. Also
\(C_F(\alpha)=\theta(C_F(c))=1\) by
Lochak--Schneps Proposition 4(i). The quotient of two
solutions is fixed by \(\alpha\), so \(h\) is unique.

Let \(Q=\exp L\) for the free class-three Lie algebra
\(L\) on two generators over \(\mathbb F_7\).
It has order \(7^5\), exponent seven, and is generated
by \(x=\exp X,y=\exp Y\); hence it is a quotient of
\(F^{(7)}\) and of \(F/R_{\mathrm{nil}}(F)\).
The third-cusp condition on \(\theta\) and A.4 give
\[
       \operatorname{coeff}_Z(\log\bar g)
                    =\frac{2^2-1}{24}=1 .
\]
But \(\psi\) has multiplier one and the same conjugator
\(g\). Peripheral preservation for \(\psi\) would force
this coefficient to be zero by A.4, a contradiction.

The exact certificate \texttt{check\_A3\_class3.py}
realizes \(Q\) in the units of
\[
 \mathbb F_7\langle X,Y\rangle/(\text{words of length }4).
\]
It enumerates all \(16807\) elements; the conjugacy
orbit of \(xy\) has \(49\) elements. For
\(g=\exp[X,Y]\) it verifies
\[
 xg^{-1}yg\not\sim xy,\qquad
 x^2g^{-1}y^2g\sim(xy)^2 .
\]
The actual Galois conjugator differs only by central
degree-three coordinates relevant to this calculation,
which do not affect it. This certifies a failure in
an actual full profinite example; no order-preserving
lift of a finite involution is being assumed.
\end{proof}

The correct universal statement from R1 is that this
\(\psi\) intertwines \(\alpha,c\) on the maximal
pro-\(p\) quotients as automorphisms of free groups.
Peripheral preservation of this specific \(\psi\)
is an additional condition and is false.

\statement{A.6}{Proved}{The pure centralizer of standard conjugation has two elements}
\[
 C_{\mathcal P}(c)=N_{\mathcal P}(\langle c\rangle)=\{1,c\}.
\]
If \(\alpha\in\widehat{GT}_0\), every conjugator in
\(\mathcal P\) from \(c\) to \(\alpha\) belongs
to \(\widehat{GT}_0\).

\begin{proof}
A normalized \((\mu,g)\) commuting with \(c\) satisfies
\(c(g)=g\), by the normalized group law. A.1 gives
\(g=1\). In particular the multiplier-one centralizer
is trivial. Let \(\Sigma=S_3\) be the geometric
permutation group. If
\(\beta c\beta^{-1}=\alpha\) and \(\alpha,c\) commute
with \(\Sigma\), then
\(r_\sigma=\beta^{-1}\sigma\beta\sigma^{-1}\)
is pure, has multiplier one, and centralizes \(c\).
Thus \(r_\sigma=1\), and the Harbater--Schneps
identification in A.1 puts \(\beta\) in
\(\widehat{GT}_0\).

Apply this to a centralizing \((\mu,1)\).
The hexagon, with \(m=(\mu-1)/2\), reads
\((xy)^m=y^mx^m\). In the finite Heisenberg groups
with multiplication
\((a,b,d)(a',b',d')=(a+a',b+b',d+d'+ab')\),
the final coordinates of these two words are
\(m(m+1)/2\) and zero. Hence \(\mu^2=1\) in
\(\widehat{\mathbb Z}\). The actual automorphism
\((\mu,1)\) then has square one. R1 forces its
multiplier to be globally \(1\) or \(-1\), proving
the centralizer assertion. A normalizer of
\(\langle c\rangle\) centralizes its unique
nonidentity element.
\end{proof}

\statement{A.7}{Proved}{The multiplier-one cohomology set already has continuum many classes}
For \(G=\mathcal P,\widehat{GT}_0\), or \(\widehat{GT}\),
put \(H_G=\ker\lambda\), with \(C_2\) acting by \(c\).
There is a canonical closed injection of profinite spaces
\[
 \widehat{\mathbb Z}^{\times}/\{1,-1\}
 \longrightarrow H^1(C_2,H_G),\qquad
 \mu\longmapsto[\theta_\mu c\theta_\mu^{-1}c],
\]
where \(\theta_\mu\) is any element of multiplier \(\mu\).
The set \(H^1(C_2,H_G)\) has cardinality \(2^{\aleph_0}\).
For \(G=\widehat{GT}_0\), \(\Delta\) vanishes on this
closed family. All these \(H_G\)-conjugacy classes form
one \(G\)-conjugacy class.

\begin{proof}
A cocycle is \(a\in H_G\) with \(a\,c(a)=1\); it
corresponds to the involution \(ac\). Cohomology
identifies \(a\) and \(h^{-1}a\,c(h)\), exactly
conjugacy by \(H_G\). The Galois image makes
\(\lambda:G\to\widehat{\mathbb Z}^{\times}\) onto.
The element \(a_\mu=\theta_\mu c\theta_\mu^{-1}c\)
is a cocycle, and replacing \(\theta_\mu\) by a
different lift conjugates its involution by \(H_G\).

Two displayed involutions are \(H_G\)-conjugate
if and only if
\(\theta_\nu^{-1}h\theta_\mu\in C_G(c)=\{1,c\}\)
for some \(h\in H_G\), by A.6. Taking multipliers
gives \(\mu=\pm\nu\); these conditions are sufficient
as well. This proves the injection.
The cocycle space is compact, and its quotient by
the compact conjugating group is Hausdorff.
The orbit map from \(G\) factors continuously
through the multiplier quotient and then its
two signs, so the injection is continuous and closed.
Its domain has cardinality continuum. Conversely
\(G\) is contained in the metrizable profinite group
\(\operatorname{Out}(F)\), of cardinality at most
continuum. This proves equality.
Finally every represented involution is conjugate
to \(c\); the character \(\Delta\) is conjugacy
invariant and \(\Delta(c)=0\).
\end{proof}

\statement{A.8}{Proved}{Full conjugacy is an additional quotient of nonabelian cohomology}
With the notation of A.7, the full \(G\)-conjugacy
classes of nonidentity involutions are
\[
 \widehat{\mathbb Z}^{\times}\backslash H^1(C_2,H_G).
\]
The action of \(\mu\), using any lift \(\theta_\mu\), is
\[
 [a]\longmapsto[\,\theta_\mu a\,c(\theta_\mu^{-1})\,].
\]
For \(\mathcal N=\ker(\mathcal P\to P_F)\), the
\(\mathcal N\)-conjugacy classes of involutions with
image exactly \(c_F\) are \(H^1(C_2,\mathcal N)\).
Full \(\mathcal P\)-conjugacy on that fibre further
quotients by \(C_{P_F}(c_F)\), acting by lifted
conjugation, which need not preserve the distinguished
cohomology class.

\begin{proof}
By R1 every nonidentity involution has multiplier
\(-1\), so is uniquely \(ac\), \(a\in H_G\).
Its square-one condition is the cocycle equation.
Conjugation by \(\theta_\mu\) gives the formula.
Different lifts differ by \(H_G\), already divided
out in cohomology, proving the first assertion.

The same calculation for \(\mathcal P/\mathcal N\)
shows that involutions over the literal \(c_F\)
are \(ac\), \(a\in\mathcal N\). A conjugation
preserves this fibre exactly when its quotient
centralizes \(c_F\). Changing a lift changes
the action only by \(\mathcal N\)-conjugacy.
Such a lift need not centralize \(c\), explaining
the last qualification.
\end{proof}

\statement{A.9}{Proved}{Odd-prime conjugacy over a relative Frattini level}
Let \(U\) be open characteristic in \(F\), let \(p\)
be odd, put \(E_p=F/O^p(U)\), \(V=U/O^p(U)\),
and consider the peripheral outer subgroup preserving
\(V\). Its kernel on
\[
           E_p/\Phi(V)=F/\overline{U'U^p}
\]
is pro-\(p\). Two involutions with the same image
in this Frattini quotient are conjugate by the kernel.
This applies to the fibre over \(c\), where that
fibre is defined by descent to \(F/U\).

\begin{proof}
\(V\) is a nonabelian finitely generated free
pro-\(p\) group. Let \(J\) be the actual
automorphisms of \(E_p\) inducing identity on
\(E_p/\Phi(V)\). Restriction maps \(J\) into
the Frattini kernel of \(\operatorname{Aut}(V)\).
That kernel is pro-\(p\): on finite lower
\(p\)-central quotients its successive kernels
embed into homomorphism groups of elementary
abelian \(p\)-groups, and those quotients are cofinal.

Restriction is injective. If an element of \(J\)
fixes \(V\) pointwise, write \(a(g)=gu_g\) with
\(u_g\in\Phi(V)\). Comparing the conjugation
actions on \(V\) gives \(u_g\in C_V(V)=1\).
Thus \(J\) is pro-\(p\). If an outer automorphism
acts inner on \(E_p/\Phi(V)\), lift its inner
conjugator to \(E_p\) and correct it to an
element of \(J\). Consequently the outer
kernel is a continuous quotient of a closed
subgroup of \(J\), hence pro-\(p\).

The common image of the two involutions has
order two, since a pro-\(p\) group for odd \(p\)
has no involution. In its inverse image the
two cyclic subgroups are complements to the
normal pro-\(p\) kernel. Finite
Schur--Zassenhaus gives conjugators in every
finite quotient; compactness gives a conjugator
in the kernel itself.
\end{proof}

For arbitrary quotient automorphisms, preservation
of the image of a merely characteristic \(U\)
must be checked. It is included in the phrase
over the class of \(c\). For fully characteristic
\(U\) it is automatic: any quotient endomorphism
lifts to an endomorphism of the free profinite
group and preserves \(U\). Such \(U\) are cofinal.

\statement{A.10}{Proved}{Every negative pure pro-\(p\) involution is locally conjugate for odd \(p\)}
Let \(P^{(p)}\) be the pure peripheral outer group
of \(F^{(p)}\), with multiplier in
\(\mathbb Z_p^\times\), for odd \(p\).
Every involution of multiplier \(-1\) is
conjugate to \(c\) by a multiplier-one element.

\begin{proof}
The multiplier kernel \(H^{(p)}\) acts trivially
on abelianization, so lies in the pro-\(p\)
Frattini kernel used in A.9. The inverse image
of \(\{1,-1\}\) under the multiplier is
\(H^{(p)}\rtimes\langle c\rangle\).
Every negative involution defines a complement;
Schur--Zassenhaus gives the required conjugacy.
\end{proof}

This does not rescue the particular map of A.5.
Its local coboundary can be changed by a local
fixed element, whereas its full coboundary is
unique. The comparison of fixed subgroups
need not be surjective.

\statement{A.11}{Proved}{A peripheral reflection triangle is equivalent to full conjugacy}
For a nontrivial involution \(\alpha\in\mathcal P\)
let \(E=F\rtimes_\alpha C_2\). Then \(\alpha\)
is conjugate to \(c\) in \(\mathcal P\) exactly
when there are involutions \(t_0,t_1,t_2\)
representing its three section classes such that
\[
 a=t_1t_0,\qquad b=t_0t_2,\qquad d=t_2t_1
\]
generate conjugates of the respective inertia
groups of \(x,y,z\), and the \(t_i\) generate
\(E\) topologically. The last condition is
equivalent to generation in every \(E/R_U\).

\begin{proof}
For standard \(c\) take
\(t_0=c,t_1=xc,t_2=y^{-1}c\); their products
are \(x,y,z\) and they generate the extension.
A peripheral conjugator transports this triple.

Conversely the universal property gives a
surjection from the profinite coproduct
\(C_2^{(0)}\amalg C_2^{(1)}\amalg C_2^{(2)}\)
to \(E\), sending the factors to the named
involutions. Its parity kernel is free
profinite on \(a,b\). The restriction maps
onto \(F\): every even word in three
reflections is a word in the two edge products.
Both kernels are free profinite of rank two,
so Hopficity makes this restriction an
isomorphism. The equality \(abd=1\) is literal.
Abelianization makes the unit exponents in
the three peripheral conditions equal,
so the induced outer map belongs to \(\mathcal P\).
Conjugation by \(t_0\) inverts both \(a\) and \(b\).
Its outer action on \(F\) is \(\alpha\), since
\(t_0\) and the named splitting element have the
same image in \(E/F\). Hence the resulting outer
map conjugates \(c\) to \(\alpha\); no choice of
an actual splitting representative is needed.

Finally put \(H=\overline{\langle t_0,t_1,t_2\rangle}\).
If \(HR_U=E\) for all \(U\), then for every \(e\in E\)
the compact sets \(H\cap eR_U\) are nonempty
and have the finite intersection property.
Their intersection and \(\bigcap R_U=1\)
imply \(e\in H\). Thus \(H=E\).
\end{proof}

The transcription of cote 145, scan 23, lists
\(l_0=t_\infty t_1,\ l_1=t_0t_\infty,\
l_\infty=t_1t_0\). With ordinary multiplication
these satisfy \(l_0l_\infty l_1=1\).
The cyclic order or path-composition convention
must therefore be made explicit before
identifying them with \(x,y,z\).
We use \(abd=1\), and do not emend an unseen scan.

\statement{A.12}{Open}{Dyadic and mixed-prime involution obstructions}
A.9 does not determine all two-primary classes or all
gluing classes over \(U=F'F^3\) and \(U_{12}\).
A.25 below constructs an actual involution with
parity one in the image of \(\widehat{GT}_0\) at
\(U_{12}\), and A.24 proves that this parity
descends to that outer image. Thus the existence
assertion at this specified relative level is proved.
Lead 3 still needs a compatible order-two full lift.
Neither R2 nor A.3 supplies that further equivariance.

For a surjection \(G\to Q\) with kernel \(L\),
a chosen lift \(t\) of an involution can be
corrected to an involution \(ut\), \(u\in L\),
exactly when
\[
                    u(tut^{-1})=t^{-2}.
\]
Surjectivity asserts the existence of \(t\),
not solvability of this nonabelian norm equation.
The map \(C_4\to C_2\) disproves that implication
in general, but is not a counterexample in
the peripheral groups. The residual task is
to solve the norm equations compatibly in
the actual relative-level kernels with the
specified frames and \(\Delta=1\), or to
prove they are insoluble.

\statement{A.13}{Open}{Parity on arbitrary involutions}
For an involution \(\tau=(-1,f)\) in
\(\widehat{GT}_0\), R3 gives
\(\Delta(\tau)=\delta((1,f))\).
It is zero for A.7's continuum family and
for every \(\widehat{GT}_0\)-conjugate of \(c\).
The cocycle relation alone gives only
\(2\Delta(a)=0\) in \(\mathbb F_2\).
A.25 gives such an involution in an actual relative
outer image of \(\widehat{GT}_0\), but no full
involution of parity one or universal full vanishing
proof is obtained, in \(\widehat{GT}_0\) or
\(\widehat{GT}\).

\statement{A.14}{Open}{The three full conjugacy assertions}
We have not decided whether every involution
of \(\mathcal P,\widehat{GT}_0,\widehat{GT}\)
is conjugate to \(c\) in its respective group.
A.2 detects conjugacy of a fixed pair on
all relative levels; A.11 identifies the
equivalent peripheral generation condition.
Neither supplies the compatible local
conjugators or reflection triangles whose
existence is required.

For \S31, folios 198--199 verso, scans
391--394, the full finite-subgroup conclusion
proved by R1 is that the only abstract
possibilities in \(\mathcal P\) are \(1,C_2\),
and \(C_2=\langle c\rangle\) occurs.
Exhaustion of its embedding conjugacy
classes is exactly the remaining assertion.

\statement{A.15}{Proved}{The pronilpotent lifting kernel also has continuum many cohomology classes}
Put \(Q=\mathcal P_F\), \(\mathcal N=\ker(\mathcal P\to Q)\),
and write \(c_Q\) for the image of \(c\). There is a canonical
closed injection
\[
 C_Q(c_Q)/\langle c_Q\rangle
 \longrightarrow H^1(C_2,\mathcal N),\qquad
 \bar\beta\longmapsto[\beta c\beta^{-1}c].
\]
Its image is exactly the \(\mathcal N\)-conjugacy classes in the
fibre over \(c_Q\) which are conjugate to \(c\) in \(\mathcal P\).
Both its domain and \(H^1(C_2,\mathcal N)\) have cardinality
\(2^{\aleph_0}\). Thus the large cohomology set does not itself
provide an additional full involution-conjugacy class.

\begin{proof}
The pronilpotent quotient is the product of its maximal
pro-\(p\) quotients. Every factor is its characteristic
Sylow subgroup; conversely, a tuple of continuous
automorphisms of the factors gives a continuous automorphism
of their product. Inner automorphisms and the three peripheral
conditions are coordinatewise. Therefore
\[
 Q=\prod_p\mathcal P^{(p)},\qquad
 C_Q(c_Q)=\prod_p C_{\mathcal P^{(p)}}(c).
\]
For odd \(p\), the multiplier map on the latter centralizer
is onto \(\mathbb Z_p^\times\). Indeed choose an element
\(\theta\) of the desired multiplier using the Galois image.
By A.10, some multiplier-one \(v\) conjugates
\(\theta c\theta^{-1}\) to \(c\); then \(v\theta\)
centralizes \(c\) and retains that multiplier.
Consequently the displayed centralizer, and its quotient
by the finite group \(\langle c_Q\rangle\), have at least
continuum many elements.

By A.2 every \(\bar\beta\in C_Q(c_Q)\) lifts to
\(\beta\in\mathcal P\). Its displayed commutator with
\(c\) lies in \(\mathcal N\) and is a cocycle.
Changing the lift changes its class only by
\(\mathcal N\)-conjugacy. Two lifts \(\beta_1,\beta_2\)
give the same class exactly when
\(\beta_2^{-1}n\beta_1\in C_{\mathcal P}(c)=\{1,c\}\)
for some \(n\in\mathcal N\). On taking the quotient,
this is exactly equality of the two cosets modulo
\(\langle c_Q\rangle\); the converse follows by lifting
that equality. This proves injectivity and the asserted
description of the image.
Compactness makes the injection continuous and closed,
as in A.7. All groups and orbit spaces here are compact
metrizable, so their cardinalities are at most continuum.
The preceding lower bound proves the asserted equalities.
\end{proof}

If the same construction is made in \(\widehat{GT}_0\),
using its actual image in any \(c\)-stable quotient,
every class in the connecting-map image has \(\Delta=0\).
This follows from conjugacy invariance of the character;
no surjectivity onto the entire peripheral quotient is
required for this observation.

\statement{A.16}{Proved}{A prescribed reflection splits every pro-two extension of the symmetric group}
Let \(1\to K\to E\xrightarrow{q}S_3\to1\) be an extension
of profinite groups with \(K\) pro-two. If a transposition
\(t\in S_3\) has an involution lift \(\widetilde t\in E\),
then there is a continuous section \(j:S_3\to E\) satisfying
\(j(t)=\widetilde t\).

\begin{proof}
For every finite \(\mathbb F_2[S_3]\)-module \(M\), restriction
\[
 H^n(S_3,M)\longrightarrow H^n(\langle t\rangle,M)
\]
is an isomorphism for \(n\geq1\). Indeed the double-coset
formula for restriction after corestriction has the identity
term and one term through the trivial intersection
\(\langle t\rangle\cap r\langle t\rangle r^{-1}=1\),
where \(r\) is a three-cycle. The second term is zero
in positive degrees. Corestriction after restriction is
multiplication by the odd index three, hence the identity
on these \(\mathbb F_2\)-vector spaces.

First suppose \(K\) is finite. Induct on its order.
When \(K\ne1\), the subgroup
\(A=\{z\in Z(K):z^2=1\}\) is a nonzero elementary
abelian characteristic subgroup of \(K\), hence normal
in \(E\). Induction gives a section of \(E/A\to S_3\)
with the prescribed image of \(\widetilde t\).
Pull back that section to an extension with kernel \(A\).
Its class in \(H^2(S_3,A)\) restricts to zero, since
\(\widetilde t\) is an involution. The proved injectivity
of restriction makes the extension split.
The difference between any splitting on \(\langle t\rangle\)
and the specified lift is an \(A\)-valued one-cocycle.
Surjectivity in degree one extends its cohomology class to
\(S_3\); changing the section by that cocycle and then by
an \(A\)-conjugation makes its value at \(t\) literally
\(\widetilde t\). This proves the finite assertion.

For general pro-two \(K\), take all open normal subgroups
\(L\) of \(E\) contained in \(K\). Each finite quotient
has such a section with the prescribed transposition.
The corresponding closed subsets of the compact space
of tuples \((e_g)_{g\in S_3}\in E^{S_3}\), imposing the
section identities modulo \(L\), are nonempty and have
the finite intersection property. Their intersection
gives a section with all identities literal.
\end{proof}

\statement{A.17}{Proved}{The parity-one relative frame can be chosen anti-invariant under conjugation}
In the notation of R2, put \(H_0=F'F^3\),
\(\Gamma_2=\widehat B/O^2(H_0)\), and let
\(d=\operatorname{Inn}(s)\circ c\), so that
\(d(s)=s\), \(d(r)=r^{-1}\), \(d(x)=y^{-1}\),
and \(d(y)=x^{-1}\). There is an actual element
\[
 b\in C_{H_0^{(2)}}(r),\qquad q(b)=[5I],\qquad d(b)=b^{-1}.
\]
Consequently the identity automorphism of
\(\widehat B/R_{\mathrm{nil}}(H_0)\), with \(s\)- and
\(\tau\)-frames one and \(r\)-frame equal to \(b\) at two
and one at all other primes, is a framed cocycle for
conjugation by \(c\). Its modulo-twelve frame parity is one.
This establishes local equivariance of these data; it does
not assert an equivariant full lift.

\begin{proof}
Use R2's actual joint quotient
\[
 J=\operatorname{im}\bigl(\widehat B\to Q\times D_0\bigr),
 \qquad Q=\mathrm{SL}_2(\mathbb Z/12)/\{\pm I\},\quad
 D_0=\widehat B/H_0.
\]
The map \(\Gamma_2\to J\) has pro-two kernel; both factors
and the joint quotient are preserved by \(d\).
The element \(z_0=([5I],1)\in J\) is central of order two
and is fixed by \(d\). Thus
\(S=\langle\bar r,z_0d\rangle\) in \(J\rtimes\langle d\rangle\)
is an actual copy of \(S_3\), with its named three-cycle
and transposition.

There is a particularly small literal lift of that
transposition as an involution. Set
\[
 C=[x,y],\qquad h=CxCx^{-1}\in F'\subseteq H_0.
\]
In R2's signed matrix coordinates modulo twelve,
\[
 q(h)=\begin{pmatrix}9&4\\8&9\end{pmatrix},\qquad
 D=\begin{pmatrix}0&1\\1&0\end{pmatrix},\qquad
 q(h)\bigl(Dq(h)D^{-1}\bigr)^{-1}=5I.
\]
Conjugation by \(D\) represents \(d\) on \(Q\).
The \(D_0\)-coordinate of \(h\) is one, so the actual
involution \(hdh^{-1}\) maps to \(z_0d\).
Apply A.16 to the preimage of \(S\) in
\(\Gamma_2\rtimes\langle d\rangle\), retaining this
involution. Its lifted three-cycle is conjugate to
the named \(r\) by the pro-two kernel, by the
coprime complement-conjugacy argument of A.9.
After that conjugation write the lifted transposition
as \(bd\). Its image in \(D_0\rtimes\langle d\rangle\)
is \(d\), so \(b\in H_0^{(2)}\), and its \(Q\)-coordinate
gives \(q(b)=[5I]\). Its order two gives \(b\,d(b)=1\).
Since it inverts the named \(r\), while \(d\) also
inverts \(r\), the element \(b\) centralizes \(r\).

All components of the prescribed relative frame have
image one in \(D_0\), so they give an actual element
of the relative fibre product. For a multiplier-one
normalized automorphism \(\Phi\) of \(\widehat B\),
write \(\Phi(s)=\beta s\beta^{-1}\) and
\(\Phi(r)=A rA^{-1}\). Direct substitution from
\(c(s)=s\), \(c(r)=sr^{-1}s\) gives the frames of
\(c\Phi c\):
\[
 \beta^c=c(\beta),\qquad
 A^c=c(\beta)s\,c(\beta)^{-1}c(A)s.
\]
For the specified local data \(\beta=1\), the second
formula is \(A^c=d(A)=A^{-1}\). These are exactly
the frames of the inverse of the identity framed
automorphism. The parity is the same as in R2's E.58:
the signed \(r\)-frame is \(5I\), whose modulo-three
value is \(-I\), and the \(s\)-frame is one.
The independently executed certificate
\texttt{check\_A17\_frame.py} verifies the literal
matrix word and the covariance identities just used.
\end{proof}

The remaining full obstruction can now be placed in a
smaller, specified kernel. Let \(L\) consist of the
multiplier-one elements whose normalized action is
identity on \(\widehat B/R_{\mathrm{nil}}(H_0)\) and whose
two pure frames both have image one there. It is the
kernel of the framed action, is stable under \(c\), and
\(\Delta(L)=0\). R2 supplies a full \(\sigma\) lifting
A.17's data, with \(\Delta(\sigma)=1\). Put \(a=\sigma c\).
Then \(a^2\in L\). A parity-one involution in this
specified fibre exists exactly when some \(v\in L\)
solves
\[
                   v\,(ava^{-1})=a^{-2}.
\]
The local construction proves that the right-hand side
lies in this framed kernel. Solvability of this last full
norm equation remains open.

\statement{A.18}{Refuted}{R2 does not admit an unrestricted order-preserving refinement}
There is an actual involution of
\(\operatorname{Out}^{\sharp}(F/R_{\mathrm{nil}}(F))\)
which lifts to a full peripheral automorphism by R2 but
has no involution lift in \(\mathcal P\).

\begin{proof}
Use the actual product
\(F/R_{\mathrm{nil}}(F)=\prod_pF^{(p)}\).
Define \(\bar a\) to be standard conjugation \(c\)
on \(F^{(3)}\) and identity on every other factor.
This is a continuous automorphism of order two,
preserving each of the three peripheral conjugacy
classes up to the common unit
\[
 \lambda_3=-1,\qquad \lambda_p=1\quad(p\ne3).
\]
Its outer class is nonidentity, since its action on
the three-primary abelianization is \(-I\).
By A.2, which applies R2 after Galois translation
of the multiplier, this class has a lift in
\(\mathcal P\). Every lift has the displayed
multiplier, because the product quotient retains
the entire abelianization \(\widehat{\mathbb Z}^2\).
By R1, a nonidentity full involution must instead
have multiplier \(-1\) at every prime. Hence no
lift of \(\bar a\) can be an involution.
\end{proof}

This counterexample is in the actual peripheral level
groups, rather than in an unrelated finite extension.
It excludes a blanket equivariant strengthening of
the central-refinement count. A.17's remaining fibre
has multiplier globally \(-1\), so this orientation
obstruction does not decide that fibre.

\statement{A.19}{Proved}{The square obstruction at an abelian refinement and the character-lifting test}
Suppose \(a\) is a lift of a quotient involution,
\(a^2\in L\), and \(L_0\triangleleft\langle L,a\rangle\)
is contained in \(L\), with \(M=L/L_0\) finite abelian.
Conjugation by \(a\) makes \(M\) a \(C_2\)-module.
An involution exists above that quotient element in
\(\langle L,a\rangle/L_0\) exactly when
\[
 [a^2]=0\quad\hbox{in}\quad
 H^2(C_2,M)=M^{a}/(1+a)M.
\]
This obstruction vanishes when \(M\) has odd order.
For elementary abelian two-primary \(M\), its target is
\(\ker(a+1)/\operatorname{im}(a+1)\).

If \(\Delta\) has a continuous lift to a character
\(\widehat{GT}_0\to\mathbb Z/4\), then it vanishes on
every involution. The same conclusion follows from
a continuous lift of \(\delta\) on \(\mathcal H\)
which is invariant under conjugation by \(c\).
Existence of either lift is not established here.

\begin{proof}
The square of the conjugation action is conjugation
by \(a^2\in L\), which is trivial on \(M\).
Also \(a^2\) is invariant under conjugation by \(a\).
Replacing \(a\) by \(va\), with \(v\in L\), changes
its square in additive notation by \((1+a)\bar v\).
This proves the quotient criterion, which is the
usual cyclic two-cocycle description of \(H^2\).
For an invariant \(u\) of odd order, \(u/2\) is
defined and \((1+a)(u/2)=u\), proving vanishing.
In characteristic two the invariants are
\(\ker(a+1)\), giving the last formula.

A character to \(\mathbb Z/4\) sends an involution
to zero or two; both reduce to zero modulo two.
Alternatively, if \(b\in\mathcal H\) is a cocycle
for \(c\) and \(\theta\) is a \(c\)-invariant
lift of \(\delta\), then
\(0=\theta(bc(b))=2\theta(b)\), so
\(\delta(b)=0\).
\end{proof}

For the actual \(a=\sigma c\) of A.17, any proposed
global character lift \(\Theta\) would necessarily satisfy
\(\Theta(a^2)=2\), although \(a^2\in L\) and
\(\Delta(L)=0\). Thus such a character must detect the
residual square inside the framed kernel; it cannot be
inferred from the local frames alone. Conversely A.17's
local parity-one framed involution rules out a character
lift to \(\mathbb Z/4\) on that local framed group itself.

\statement{A.20}{Proved}{The locally glued symmetric subgroup does not lift in the standard extension}
Let \(H_0\), \(d\), and \(b\) be as in A.17. In the standard
relative extension
\[
 \bigl(\widehat B/R_{\mathrm{nil}}(H_0)\bigr)
                  \rtimes\langle d\rangle
\]
there is a subgroup \(S_{\mathrm{loc}}\simeq S_3\), generated
by the named \(r\) and a transposition with component \(bd\)
at two and \(d\) at every odd prime. No finite subgroup
of the standard full extension
\(\widehat B\rtimes\langle d\rangle\) maps onto
\(S_{\mathrm{loc}}\). The obstruction is visible already in
\(Q\rtimes\langle d\rangle\), with \(|Q|=576\).

\begin{proof}
All local components of the transposition have the same
image \(d\) in \(D_0\rtimes\langle d\rangle\).
They therefore form an actual element of the relative fibre
product. A.17 gives its order two and its inversion of \(r\).
Their images in \(D_0\rtimes\langle d\rangle\) generate
\(S_3\), so the resulting subgroup has order exactly six.

In the full free product \(\widehat B=C_2\amalg C_3\),
every element of order three is conjugate into the named
\(C_3\), and that factor is malnormal. These are
Herfort--Ribes Theorems A' and B'.\cite{A_HR1985}
In particular \(N_{\widehat B}(\langle r\rangle)=\langle r\rangle\).
Since \(d\) inverts \(r\),
\[
 N_{\widehat B\rtimes\langle d\rangle}(\langle r\rangle)
                         =\langle r,d\rangle\simeq S_3.
\]
Every full subgroup isomorphic to \(S_3\) intersects
\(\widehat B\) in its unique subgroup of order three:
it cannot be contained in \(\widehat B\), whose finite
subgroups conjugate into its factors by the same paper's
Theorem 1. Conjugating that order-three subgroup into
\(\langle r\rangle\) proves that every full \(S_3\)
is conjugate to \(\langle r,d\rangle\).

The image of \(S_{\mathrm{loc}}\) in the actual finite
quotient \(Q\rtimes\langle d\rangle\) is
\[
                   \langle R,[5I]d\rangle .
\]
It is not conjugate there to \(\langle R,d\rangle\).
Indeed a conjugator may be taken in \(Q\), since \(d\)
normalizes the latter subgroup, and it must normalize
its unique order-three subgroup. Directly solving the
two-by-two equations in \(\mathrm{SL}_2(\mathbb F_3)\)
shows that no determinant-one matrix conjugates \(R\)
to \(R^{-1}\), and its centralizer is \(\{\pm R^j\}\).
The centralizer modulo four is also \(\{\pm R^j\}\).
Traces exclude a negative sign in the projective
commutation equation. Thus a representative of a
normalizing element has coordinates
\(\varepsilon_4R^i\) modulo four and
\(\varepsilon_3R^j\) modulo three.
Conjugation by \(D\) fixes both signs and inverts \(R\).
Its \(d\)-norm consequently has positive signs in both
coordinates. It cannot lie in \([5I]\langle R\rangle\):
one global choice of representative has negative sign
modulo three, and the other has negative sign modulo four.
This proves the nonconjugacy and hence the asserted
nonlifting of an \(S_3\) subgroup.

For completeness, a finite subgroup mapping onto
\(S_{\mathrm{loc}}\) has trivial kernel: the kernel lies
in the torsion-free free profinite group \(H_0\).
It would therefore be an \(S_3\), already excluded.
The exact exhaustive certificate
\texttt{check\_A20\_standard\_extension.py} generates all
\(576\) elements of \(Q\), obtains normalizer and
centralizer orders \(18\), lists all nine \(d\)-norm
values and all three elements of the target coset, and
checks that the lists are disjoint. It independently
enumerates the two signed centralizers used above.
\end{proof}

A.20 blocks lifting this subgroup inside the \emph{standard}
extension. A hypothetical new involution of
\(\widehat{GT}_0\) would change the action defining the
full extension. No obstruction in that unknown twisted
extension is inferred from A.20.

\statement{A.21}{Refuted}{Real projectivity and a rank-two open free subgroup do not imply finite graph recognition}
There exists a finitely generated real-projective profinite
group \(E_*\) with an open normal subgroup isomorphic to
\(\widehat F_2\), quotient \(C_2\), and self-centralizing
involutions, which is not the profinite fundamental group
of a finite graph of finite groups. It has one conjugacy
class of involutions, so it does not satisfy R1's
three-section conclusion.

\begin{proof}
Put \(A=\prod_{p\text{ odd}}\mathbb Z_p\) and define
\[
 \Gamma_*=A\rtimes\mathbb Z_2,
 \qquad u\cdot a=(-1)^{u\bmod2}a.
\]
All groups and actions here are profinite and continuous.
The kernel of the parity map \(\Gamma_*\to C_2\) is
\(A\times2\mathbb Z_2\simeq\widehat{\mathbb Z}\).
The group \(\Gamma_*\) is torsion-free: a finite-order
element projects trivially to \(\mathbb Z_2\), and then
belongs to the torsion-free group \(A\).
Set
\[
 E_*=C_2\amalg\Gamma_*,\qquad
 \varepsilon:E_*\longrightarrow C_2,
\]
where \(\varepsilon\) is identity on the first factor
and parity on the second. In the open-subgroup Kurosh
decomposition of \(\ker\varepsilon\), the two vertex
groups are \(1\) and \(\widehat{\mathbb Z}\), and the
quotient graph has two vertices joined by two edges.
Consequently
\[
 \ker\varepsilon\simeq
 \widehat{\mathbb Z}\amalg\widehat{\mathbb Z}
 \simeq\widehat F_2.
\]
The open Kurosh theorem in precisely this full-profinite
form is recalled in Herfort--Ribes, Section~2.
Their Theorems A', B', and 1 show that all nonidentity
torsion lies in conjugates of the single \(C_2\)
factor, and that each such involution has centralizer
exactly that factor.\cite{A_HR1985}
Haran's characterization of real projectivity therefore
applies: an index-at-most-two subgroup of cohomological
dimension at most one, together with self-centralizing
involutions, suffices. We use its exact formulation in
P\'al--Quick, Theorem~2.1.\cite{A_PalQuickFormal2025}

Suppose \(E_*\) were the profinite fundamental group of
a finite graph of finite groups. Its vertex groups have
order at most two. Collapse each nonloop edge whose
nontrivial edge group is equal to an endpoint group.
In the resulting reduced graph, a loop with edge group
\(C_2\) would make its stable letter centralize that
\(C_2\); the stable letter has infinite order, as is
seen by killing the vertex groups and mapping that
loop to \(\widehat{\mathbb Z}\). This contradicts the
centralizer just proved. All remaining edges therefore
have trivial group. We obtain
\[
 E_*\simeq\bigl(\mathop{\amalg}_{j=1}^{k}C_2\bigr)
                   \amalg\widehat F_r .
\]
The single torsion class forces \(k=1\). The
torsion-free index-two subgroup of this free product
has rank \(2r+k-1\), by its two-sheeted graph or by
open Kurosh. The rank two already computed forces
\(r=1\).

On the other hand \([\Gamma_*,\Gamma_*]=A\), since
multiplication by two is invertible on \(A\).
Thus
\[
 E_*^{\mathrm{ab}}=C_2\times\mathbb Z_2,
\]
whereas \((C_2\amalg\widehat{\mathbb Z})^{\mathrm{ab}}
=C_2\times\widehat{\mathbb Z}\). The latter has
nontrivial odd-primary quotients and the former does not.
This contradiction proves the claim.
\end{proof}

The construction uses the torsion-free virtually
procyclic mechanism already present in Haran's
counterexamples to full-profinite virtual freeness
recognition.\cite{A_Haran1985}
It identifies a genuine limitation of that method;
it is not a new peripheral involution. For the groups
in A.14 the three-section hypothesis, and ultimately
the marked cusp products in A.11, still have to be used.

\statement{A.22}{Proved}{The three distinguished reflections normally generate, and a proper generated subgroup has infinite index}
Let \(\alpha\) and \(h\) be as in R1, and put
\(E=F\rtimes_\alpha\langle t\rangle\), \(b=hyh^{-1}\).
The three involutions
\[
       t_0=t,\qquad t_1=xt,\qquad t_2=b^{-1}t
\]
represent the three section classes and normally
generate \(E\). If their closed generated subgroup
\(J\) is proper, then \([E:J]\) is infinite.
Normal generation alone does not supply the full
generation required in A.11.

\begin{proof}
The equality \(f=h^{-1}\alpha(h)\) gives
\(\alpha(b)=b^{-1}\), so all three displayed
elements have order two. In
\(E^{\mathrm{ab}}\simeq C_2^3\), their images are
\(t,x+t,y+t\); these are distinct and hence the
elements are pairwise nonconjugate. R1 states that
there are exactly three classes. Quotienting \(E\)
by their closed normal closure kills \(t\), then
\(x=t_1t\), then \(b^{-1}=t_2t\). Since \(b\) is a
conjugate of \(y\), it kills \(y\), and the quotient
is trivial.

There is a continuous epimorphism
\(C_2\amalg C_2\amalg C_2\to J\) sending the three
factors to the displayed involutions. Its parity
kernel is free profinite of rank two, so \(J\cap F\)
is generated by at most two elements. If \(J\)
had finite index in \(E\), then \(J\cap F\) would
be open in \(F\) and the profinite Schreier formula
would give
\[
 d(J\cap F)=1+[F:J\cap F].
\]
It follows that \(J\cap F=F\), and \(t\in J\)
then gives \(J=E\). This proves the assertion.
\end{proof}

\statement{A.23}{Proved}{Torsion and cusp centralizers in the relative primary modular groups}
Let \(q:\widehat B=C_2\amalg C_3\to D\) be a finite
quotient injective on the two factors, with torsion-free
kernel \(U\), and write
\[
 \Gamma_p=\widehat B/O^p(U),\qquad V_p=U^{(p)}.
\]
For \(a=s\) or \(r\), of prime order \(\ell\), the
following statements hold:
\[
 \begin{array}{ll}
 p\ne\ell:&C_{\Gamma_p}(a)\longrightarrow C_D(q(a))
                    \text{ is onto},\\[2pt]
 p=\ell:&C_{\Gamma_p}(a)=\langle a\rangle.
 \end{array}
\]
For \(U=U_{12}=\ker q_{12}\) in the projective modular
quotient, put \(\Gamma=\widehat B/R_{\mathrm{nil}}(U)\).
Then
\[
 C_{\Gamma_p}(\tau)=\overline{\langle\tau\rangle},
 \qquad C_\Gamma(\tau)=\overline{\langle\tau\rangle}
                    \simeq\widehat{\mathbb Z}.
\]
In the second formula the closure is taken in \(\Gamma\);
in the first it is taken in \(\Gamma_p\).

\begin{proof}
If \(p\ne\ell\), choose a lift \(g\in\Gamma_p\)
of an element of \(C_D(q(a))\). The subgroups
\(\langle a\rangle\) and
\(\langle gag^{-1}\rangle\) are complements to
\(V_p\) in the inverse image of \(\langle q(a)\rangle\).
Coprime Schur--Zassenhaus, first in finite quotients
and then by compactness, conjugates the second to
the first by an element of \(V_p\). The generators
with image \(q(a)\) then agree. Adjusting \(g\)
by that element gives the asserted centralizer lift.

Now let \(p=\ell\), and put
\(J=q^{-1}(\langle q(a)\rangle)\subset\widehat B\).
The open Kurosh decomposition expresses \(J\) as
a finite free profinite product of copies of
\(C_p\), together with a finitely generated free
profinite factor. It has no factor of the other
prime order. Since \(J/U=C_p\),
\[
          J/O^p(U)=J^{(p)}.
\]
Indeed the left side is pro-\(p\), and every
homomorphism from \(J\) to a finite \(p\)-group
kills \(O^p(U)\). Its maximal pro-\(p\) quotient
is the corresponding free pro-\(p\) product.

If \(g\in C_{\Gamma_p}(a)\), choose
\(g_0\in\widehat B\) with the same \(D\)-image.
The elements \(a\) and \(g_0ag_0^{-1}\) belong
to \(J\) and become conjugate by \(V_p\) in
\(J^{(p)}\). Distinct Kurosh factors remain
nonconjugate in that free pro-\(p\) product,
by Herfort--Ribes Theorem B'. Their double-coset
labels in \(J\backslash\widehat B/\langle a\rangle\)
must consequently agree. Thus \(g_0\in J\),
so \(g\in J^{(p)}\). The same theorem gives
\(C_{J^{(p)}}(a)=\langle a\rangle\).
This proves the prime-primary assertion without
assuming a full-profinite Kurosh theorem for
arbitrary closed subgroups.\cite{A_HR1985}

For the cusp assertion, \(|D|=576\) and
\(q(\tau)\) has order twelve. The discrete kernel
has 48 cusps, genus 25, and free rank 97:
the cusp count is \(576/12\), and its Euler
characteristic is \(576(-1/6)=-96\).
In \(H_1(U,\mathbb F_p)\), the 48 oriented cusp
vectors have precisely the relation that their
sum is zero. In particular they are pairwise
distinct for every prime, including two and three.
The \(D\)-stabilizer of the vector represented
by \(\tau^{12}\) is exactly
\(\langle q(\tau)\rangle\): deck transformations
permute these vectors by their cusp permutation.
These assertions follow directly from the usual
oriented surface presentation of this torsion-free
finite cover of the modular orbifold.

If \(g\in\Gamma_p\) centralizes \(\tau\), it
centralizes \(\tau^{12}\in V_p\). Its image in
\(D\) therefore belongs to the indicated cusp
stabilizer. After multiplying \(g\) by an integral
power of \(\tau\), it belongs to \(V_p\).
The boundary element \(\tau^{12}\) is primitive
in \(U\), since another cusp can be omitted in
the surface free basis, and remains primitive
in \(V_p\). Herfort--Ribes Theorem B' in that
free pro-\(p\) group gives
\(C_{V_p}(\tau^{12})=\overline{\langle\tau^{12}\rangle}\).
This proves the first cusp equality. The fibre
product of these procyclic groups over
\(\langle q(\tau)\rangle=C_{12}\) is exactly
\(\widehat{\mathbb Z}\): its \(p\)-coordinate
is the infinite \(p\)-primary exponent, and the
remaining finite coordinates are the common
residue modulo twelve. This proves the second.
The same proof applies to \(x=\tau^{-2}\): its
negative sixth power is \(\tau^{12}\), so a
centralizer of \(x\) fixes the same cusp vector.
Conjugating by \(s\) treats \(y=sxs\). Thus
\(C_\Gamma(x)=\overline{\langle\tau\rangle}\)
and \(C_\Gamma(y)=s\overline{\langle\tau\rangle}s\),
as used below.
\end{proof}

\statement{A.24}{Proved}{The parity character descends to the actual outer image at level twelve}
Let
\[
 F_{12}=F/R_{\mathrm{nil}}(U_{12}),\qquad
 G_{12}=\operatorname{im}\!\left(
 \widehat{GT}_0\longrightarrow
               \operatorname{Out}^{\sharp}(F_{12})\right).
\]
The map is well-defined by the finite covariance in
refereed LM7\_PUSH E.65, which preserves \(U_{12}\).
There is a unique continuous character
\(\Delta_{12}:G_{12}\to\mathbb F_2\) such that
\(\Delta=\Delta_{12}\circ\pi_{12}\).

\begin{proof}
First consider a normalized multiplier-one action
\(\Phi\) which is identity on
\(\Gamma=\widehat B/R_{\mathrm{nil}}(U_{12})\).
Its pure \(s\)-frame, when projected to \(\Gamma_2\),
centralizes \(s\), so A.23 puts it in
\(\langle s\rangle\). Purity makes this projection
identity, since \(F\cap\langle s\rangle=1\)
already in the quotient \(S_3\). Hence its
\(Q_{12}\)-image is identity. At \(p=3\) the
same argument applies to the pure \(r\)-frame.
Moreover \(\Phi(s)=sf\) gives \(q_{12}(f)=1\).
The explicit covariance formula of LM7\_PUSH
E.65.9 then gives \(\rho_\Phi\equiv0\pmod3\),
and E.65.11 gives \(\rho_\Phi\in2\widehat{\mathbb Z}\).
Thus \(V_\Phi=I\) modulo twelve. Both parity
bits are zero. This proves that \(\Delta\)
vanishes on the kernel of the actual \(B\)-action.

We check that passing to an outer action on
\(F_{12}\) adds no kernel. Its abelianization is
\(\widehat{\mathbb Z}^2\), because
\(R_{\mathrm{nil}}(U_{12})\subset U_{12}'\subset F'\).
An outer-trivial action therefore has multiplier
one. Suppose its restriction is conjugation by
\(u\in F_{12}\). The centralizer of \(F_{12}\)
in \(\Gamma\) is trivial. To see this, the cusp
argument in A.23 also gives
\(C_\Gamma(x)=\overline{\langle\tau\rangle}\)
and
\(C_\Gamma(y)=s\overline{\langle\tau\rangle}s\).
Their images in \(S_3\) are distinct order-two
subgroups; an element of their intersection lies
in \(F_{12}\), and then in both
\(\overline{\langle x\rangle}\) and
\(\overline{\langle y\rangle}\). Abelianization
shows that their intersection is trivial.

An automorphism of \(\Gamma\) fixing the open
normal subgroup \(F_{12}\) pointwise is consequently
identity: for each \(g\in\Gamma\), its discrepancy
\(\Phi(g)g^{-1}\) centralizes that subgroup.
Thus the given \(\Phi\) is conjugation by \(u\)
on all of \(\Gamma\). It fixes \(\tau\), so A.23
gives \(u\in\overline{\langle\tau\rangle}\cap F_{12}
=\overline{\langle x\rangle}\). Write \(u=x^a\).
The normalized formula \(\Phi(s)=sf\) then gives
\(f=y^a x^{-a}\). Since \(f\in F_{12}'\), its
abelian coordinates \((-a,a)\) vanish; hence
\(a=0\). The actual action was identity after all.

We have proved kernel containment for \(\Delta\).
It therefore factors through \(G_{12}\); continuity
follows because a continuous surjection of compact
groups to a Hausdorff group is a quotient map.
\end{proof}

\statement{A.25}{Proved}{An actual parity-one involution exists at the relative level twelve}
The group \(G_{12}\) of A.24 contains an involution
\(a_{12}\) with
\[
             \lambda(a_{12})=-1,
             \qquad\Delta_{12}(a_{12})=1.
\]
It is not conjugate to the image of \(c\) in
\(G_{12}\). Thus the first, specified-relative-level
assertion of Lead~3 holds at \(U_{12}\).
The existence of an involution lifting it in the
full group is not asserted.

\begin{proof}
Let \(r_0=(3/4)^{1/4}>0\). The splitting field
\(L=\mathbb Q(i,r_0)\) has degree eight: the
polynomial \(4X^4-3\) is Eisenstein at three,
and \(\mathbb Q(r_0)\) is real. Choose a Galois
element \(\gamma\) whose restriction to \(L\)
fixes \(i\) and sends \(r_0\) to \(ir_0\).
Such a restriction is an automorphism of \(L\)
and extends to \(\overline{\mathbb Q}\).
Use its actual Galois image in \(\widehat{GT}\),
and put
\[
 \sigma_0=\gamma c\gamma^{-1}c,
 \qquad \sigma=\sigma_0^3.
\]
Then \(\lambda(\sigma_0)=1\) and
\(c\sigma_0c=\sigma_0^{-1}\), so
\(\sigma c\) is an actual full involution.
Also \(\sigma_0(r_0)=-r_0\). The Galois comparison
in refereed LM7\_PUSH E.65.13, including the model
and sign calculation supplied by its referee,
makes both normalized CM bits of \(\sigma_0\) equal
to one. More precisely these are
\(\varepsilon_3(T_{\sigma_0})\) and
\(\varepsilon_2(W_{\sigma_0})\), with the shear
\(V\) included in the definitions of \(T,W\).
LM7\_PUSH E.65.12 proves that \(T,W\) are
homomorphisms on \(\mathcal H\); composing with
the two indicated characters proves that both
bits are characters on all of \(\mathcal H\).
The same bits therefore hold for \(\sigma\).

The cubing calculation in the last paragraph of
that refereed proof gives
\[
 q\Phi_\sigma=q,
 \qquad q_{12}(f_\sigma)=I,
 \qquad V_\sigma=I.
\]
Let \(A,\beta\in F\) be its actual pure frames,
and write \(a=q(A),\ b=q(\beta)\). The r-bit
and s-bit are both one, while \(a\) and \(b\)
centralize \(q(r)\) and \(q(s)\), respectively.
More precisely, cubing the two finite cyclic
characters gives
\[
 a=[5I],\qquad
 b=[B]\ \text{or}\ [B^{-1}],\qquad
 B=\begin{pmatrix}9&4\\8&9\end{pmatrix}\pmod{12}.
\]
Here \(B\equiv I\pmod4\), \(B\equiv S\pmod3\),
and \(B^2=5I\). The choice between the two
possibilities will not matter.

On \(\Gamma_p=\widehat B/O^p(U_{12})\), define
\[
 \phi_2=\Phi_\sigma,\qquad
 \phi_p=\operatorname{id}\quad(p\ne2).
\]
All induce identity on the shared finite quotient
\(Q_{12}\). They therefore define an actual
continuous automorphism \(\phi\) of their fibre
product \(\Gamma\). Each fixes \(\tau\).
The automorphism \(\phi c\) has square one:
at two it is \(\Phi_\sigma c\), and at every
other prime it is \(c\).

We now verify full peripheral gluing with literal
conjugators, the essential point of the construction.
At two, A.23 supplies \(k_2\in C_{\Gamma_2}(r)\)
with \(q(k_2)=a^{-1}\). Replace the r-frame
\(A_2\) by \(A_2k_2\). Its \(Q_{12}\)-image
is one. At every other prime take the identity
r-frame. These give a single pure r-frame of
\(\phi\), with \(Q_{12}\)-image one.

At two retain the s-frame \(\beta_2\), with
\(Q_{12}\)-image \(b\). For every odd prime,
A.23 supplies \(k_p\in C_{\Gamma_p}(s)\)
with \(q(k_p)=b\). These are frames for the
identity automorphism at those primes. Their
common image \(b\in q(F)\) makes the resulting
global s-frame pure. The \(\tau\)-frame is one
at every prime. Thus \(\phi\) preserves all
three conjugacy classes \(s,r,\tau\), with
the stated common quotient frames.

Apply R2 with these literal conjugators. It gives
an automorphism of \(\widehat B\) fixing \(\tau\)
and lifting \(\phi\), with pure frames whose
\(Q_{12}\)-images are \((1,b)\) in the order
\((r,s)\). This lift belongs to \(\mathcal H\).
Indeed write it as \(s\mapsto sf\),
\(r\mapsto f^{-1}r\). Its reduction has
\(f_2=f_\sigma\) and \(f_p=1\) at odd primes,
so its image in every primary abelianization of
\(F\) is zero. Since
\(R_{\mathrm{nil}}(U_{12})\subset F'\), the
full \(f\) lies in \(F'\). The two torsion
relations are exactly the gentle relations, and
R2 supplies invertibility.

Denote this full multiplier-one lift by
\(\widetilde\sigma\). Its action on \(Q_{12}\)
is identity and its \(f\)-image there is one;
the covariance formula makes \(V=I\). Its
r-bit is zero and its s-bit is one. Hence
\(\Delta(\widetilde\sigma)=1\). The outer
image
\[
                  a_{12}=\pi_{12}(\widetilde\sigma c)
\]
is precisely the involution already constructed.
Its action on \(F_{12}^{\mathrm{ab}}\) is \(-I\),
so it has order exactly two. A.24 gives
\(\Delta_{12}(a_{12})=1\), whereas
\(\Delta_{12}(\pi_{12}(c))=0\). The character
separates their conjugacy classes.
\end{proof}

The exact finite computations underlying the construction
are retained in \path{check_A25_level.py} and its output:
the quotient has order 576, the pure subgroup has order 96,
the cusp orbit has size 48, and the frame matrix has
square \([5I]\). The script also verifies the dihedral
action on the four roots used above. The centralizer
and lifting assertions are proved in the text, not
inferred from this finite computation.

Unlike A.17, A.25 is an actual involution in an
outer image of the relative level. It uses the
different primes where torsion centralizers acquire
extra lifts, and retains their common finite quotient
frames throughout. Let
\(L_{12}=\ker(\widehat{GT}_0\to G_{12})\) and
\(a=\widetilde\sigma c\). Then
\(a^2\in L_{12}\) and \(\Delta(L_{12})=0\).
A full involution above this specified class exists
if and only if
\[
          v(ava^{-1})=a^{-2}
          \quad\text{has a solution }v\in L_{12}.
\]
A.19 describes every abelian refinement of this
remaining square obstruction. No solution or
incompatible finite refinement has yet been proved.

\statement{A.26}{Proved}{Finite frame visibility prevents arithmetic-only continuation}
Put
\[
 U=U_{12},\qquad V=\overline{U'U^6},\qquad D_6=\widehat B/V.
\]
This is a finite quotient of order \(576\cdot6^{97}\),
and \(V\) is characteristic in \(U\), normal in
\(\widehat B\), and preserved by \(\widehat{GT}_0\).
The two normalized CM bits of an element of
\(\mathcal H\) are determined by its actual action
on \(D_6\). In particular both bits are zero for
an element acting identically on \(D_6\).

Let \(H_6\) be the actual image of \(\mathcal H\)
in \(\operatorname{Aut}(D_6)\), and
\(J_6=\langle H_6,c_6\rangle\). There is a
character \(\Delta_6:J_6\to\mathbb F_2\) induced
by \(\Delta\) on
\(\mathcal H\rtimes\langle c\rangle\).
The involution of A.25 induces an actual involution
\(a_6\in J_6\) with \(\Delta_6(a_6)=1\).
No Galois element of multiplier \(-1\) has the
same actual action on \(D_6\).

\begin{proof}
For \(p=2,3\), write
\(D_p=\widehat B/\overline{U'U^p}\).
If \(a=s\) when \(p=2\), or \(a=r\) when
\(p=3\), then
\[
  \operatorname{im}\bigl(C_{D_p}(a)\to Q_{12}\bigr)
                      =\langle q(a)\rangle.
\]
To prove this, let \(g\) be a centralizing element
and choose a lift \(g_0\in\widehat B\) of its
\(Q_{12}\)-image. In
\(J=q^{-1}(\langle q(a)\rangle)\), the elements
\(a\) and \(g_0ag_0^{-1}\) become conjugate
by an element of \(U/\overline{U'U^p}\).
If their Kurosh factor labels were different,
the free-product universal property would give
a continuous character \(J\to C_p\) sending
one to one and the other to zero. Every such
character kills \(\overline{U'U^p}\), and
characters are constant on conjugacy classes.
This is a contradiction. Therefore the labels
coincide and \(q(g_0)\in\langle q(a)\rangle\),
as claimed. The reverse inclusion is immediate.

Any two pure s-frames for the same actual action
on \(D_6\) differ by a centralizer element.
The displayed assertion at two, followed by
intersection with the pure subgroup, shows that
their \(Q_{12}\)-images agree. The assertion at
three does the same for r-frames. The shear
\(V_\Phi\) modulo twelve is determined by the
\(Q_{12}\)-image of \(f\), by the explicit
covariance formula and the evenness of \(\rho\).
Thus both normalized bits are determined.

The kernel of \(\mathcal H\to H_6\) therefore
lies in \(\ker\delta\), giving a character
\(\delta_6\) on \(H_6\). It is invariant under
\(c_6\) because \(\delta\) is c-invariant.
The element \(c_6\) is not in \(H_6\):
on the retained abelian quotient
\(\widehat B^{\mathrm{ab}}=C_6\), its action
is inversion, whereas \(H_6\) acts identically.
Hence \(J_6=H_6\rtimes\langle c_6\rangle\),
and
\(\Delta_6(hc_6^e)=\delta_6(h)\) is well-defined.

The group \(U/R_{\mathrm{nil}}(U)\) has
characteristic subgroup whose quotient is
\(U/\overline{U'U^6}\). The actual involution
\(\phi c\) in A.25 therefore induces an actual
involution on \(D_6\). It belongs to \(J_6\),
has \(\Delta_6=1\), and is nonidentity since
it inverts the abelian quotient \(C_6\).
If a Galois element of multiplier \(-1\)
induced this action, it would belong to
\(\mathcal H\rtimes\langle c\rangle\) and
its \(\Delta\)-value would be one. This
contradicts R3's Galois vanishing.
\end{proof}

This is an obstruction to continuing the
\emph{arithmetic-only construction}, not an
obstruction to lifting \(a_{12}\) at all.
At any relative level below \(V\), the common
finite action must be \(a_6\). Thus no primary
component of such a continuation can be obtained
from a full Galois element of multiplier \(-1\).
In particular one cannot choose identity-on-
\(D_6\) Galois c-cocycles with the two CM bits
still equal to one. Nongalois relative lifts,
or a solution of the norm equation in
\(L_{12}\), remain possible. This conclusion
does not identify the entire character
\(\Delta\) on the image of arbitrary multipliers
in \(\operatorname{Aut}(D_6)\); the precise
finite domain here is \(J_6\).

\statement{A.27}{Proved}{The exact finite real-reflection count and its properness condition}
Let \(\pi:E\twoheadrightarrow D\) be an epimorphism of
finite groups. Fix involutions \(d_0,d_1,d_2\in D\)
which generate \(D\), and subsets
\(T_{ij}\subseteq\pi^{-1}(d_id_j)\), for
\((i,j)=(0,1),(1,2),(2,0)\). For every subgroup
\(J\leq E\), put
\[
 I_i(J)=\{u\in J:\pi(u)=d_i,\ u^2=1\},
 \qquad
 M_{ij}(J)_{u,v}=\mathbf1_{T_{ij}}(uv).
\]
Thus \(M_{ij}(J)\) is a rectangular matrix with rows
indexed by \(I_i(J)\) and columns indexed by
\(I_j(J)\). The number of triples of involutions
in \(E\) with the specified images and pair products,
and which generate \(E\), is
\[
 \sum_{J\leq E}\mu(J,E)
       \operatorname{tr}\bigl(M_{01}(J)M_{12}(J)M_{20}(J)\bigr),
 \tag{A.27.1}
\]
where \(\mu\) is the ordinary subgroup-lattice
M\"obius function. Empty index sets contribute zero.
This count concerns actual homomorphisms from
\(C_2*C_2*C_2\), with actual quotient data.

\begin{proof}
Expanding the trace gives exactly one summand for
\((u_0,u_1,u_2)\in I_0(J)\times I_1(J)\times I_2(J)\),
and that summand is one exactly when all three
pair-product conditions hold. Their ordered product
is automatically one:
\((u_0u_1)(u_1u_2)(u_2u_0)=1\).
Every counted triple has a unique generated subgroup.
M\"obius inversion on the finite subgroup lattice
therefore gives (A.27.1). The assumption that the
\(d_i\) generate makes every subgroup not mapping
onto \(D\) contribute zero.
\end{proof}

The matrices in (A.27.1) retain the common reflection
at the ends of adjacent edges. The ordinary
boundary-class trace sums of R2 do not retain this
incidence datum. Equality of those ordinary sums
therefore does not by itself establish equality of
(A.27.1). A concrete perfect-kernel control is
\(S_5\twoheadrightarrow C_2\), with kernel \(A_5\),
and all three \(d_i\) nonidentity. Every \(d_i\)
has an involution lift, and \(S_5\) has three
ordinary generators, but three such lifts cannot
generate \(S_5\): they are transpositions, whose
three edges cannot connect a graph on five vertices.
The certificate \path{check_A27_real_count.py}
checks perfectness of the kernel and enumerates
all 1000 triples. The generated-subgroup orders
and multiplicities are
\[
 (2,10),\quad(4,90),\quad(6,240),\quad
 (12,180),\quad(24,480).
\]
This is a control on a proposed general real
embedding theorem, not a counterexample to the
negative pants problem. Indeed a product of two
odd involutions in \(S_5\) has order at most three;
no generating pants triple of \(A_5\) has all
three orders at most three. The same certificate
checks all 816 candidate ordered pairs and finds
no pair generating \(A_5\).

\statement{A.28}{Proved}{A negative finite peripheral action is stopped by a central spin obstruction}
There is an epimorphism \(q:F\twoheadrightarrow A_6\)
and an involutory automorphism \(\vartheta\) of
\(A_6\) which sends each of the three labelled
cusp images to an \(A_6\)-conjugate of its inverse,
but no multiplier-minus-one peripheral automorphism
of \(F\) satisfies
\[
                       q\Phi=\vartheta q.
\]
In particular no full negative peripheral involution
induces this finite action. The obstruction is
already visible in a central double cover; it is
not a counterexample to R2 or A.3.

\begin{proof}
We use rightmost-first permutation composition and put
\[
\begin{aligned}
 X_0&=(1\,2\,3\,4)(5\,6),&
 Y_0&=(1\,2\,4\,6)(3\,5),\\
 Z_0&=(1\,2\,6\,3)(4\,5),&
 t&=(2\,4).
\end{aligned}
\]
The certificate \path{check_A28_s6_passport.py}
verifies \(X_0Y_0Z_0=1\) and
\(\langle X_0,Y_0\rangle=A_6\), giving
\(q(x)=X_0,q(y)=Y_0,q(z)=Z_0\).
Let \(\vartheta=\operatorname{Ad}(t)\).
The required inverse-class conjugators are,
in the displayed cusp order,
\[
                    1,\qquad(2\,6\,4),\qquad(2\,3\,4).
\]
In particular these conjugators are even.
Each cusp also has an odd involution which
inverts it: respectively \((2\,4),(2\,6),(2\,3)\).
The permutation \((3\,4\,6)\), of order three
in \(A_6\), cyclically permutes the three
cusp images by conjugation.

For completeness there is a short section-incidence
explanation of why an order-two full lift is impossible.
Odd involutions of \(S_6\) have just two cycle types,
\((2)\) and \((2,2,2)\). The product of two of
one type cannot have cycle type \((4,2)\).
For transpositions its order is at most three.
For two fixed-point-free involutions, their two
perfect matchings decompose into alternating cycles;
on six letters the possible product orders are
again at most three. Thus every \((4,2)\) cusp
requires opposite types at its two reflection ends.

In a hypothetical full negative pants extension,
R1 gives exactly three section classes. Their
images in
\(E^{\mathrm{ab}}=C_2^3\) are
\(\bar t,\bar x+\bar t,\bar y+\bar t\).
The x-cusp joins the first two, the y-cusp the
first and third, and the z-cusp the last two:
its two ends differ by \(\bar x+\bar y\).
Here each cusp has genuine involution ends.
Indeed an element of the odd coset which inverts
a full cusp has square in its procyclic centralizer;
inversion fixes that square, so the square is one.
Their three finite class labels would therefore
colour a triangle with two colours, with opposite
colours on each edge. This is impossible. No
assumption that the three reflections generate
is made in this argument.

A stronger obstruction excludes even infinite-order
full peripheral lifts of multiplier \(-1\).
We give a finite certificate rather than assume
that an arbitrary finite peripheral automorphism
lifts. The two degree-80 permutations in the
ATLAS representation \texttt{2S6G1-p80B0}
are opened source data. The code
\path{check_A28_double_cover.py} enumerates their
joint group with the natural standard generators
\((1\,2),(2\,3\,4\,5\,6)\) of \(S_6\).
It verifies an epimorphism
\[
 \widetilde S\twoheadrightarrow S_6,
 \qquad |\widetilde S|=1440,
 \qquad \ker=\langle\epsilon\rangle\cong C_2,
\]
with central kernel. Its even preimage
\(\widetilde A\) has order 720. Choose arbitrary
lifts \(X,Y\in\widetilde A\) of \(X_0,Y_0\),
and put \(Z=(XY)^{-1}\). Choose an involution
\(T\in\widetilde S\) above \(t\).
The enumeration of all conjugators in
\(\widetilde A\) proves, for each
\(Q\in\{X,Y,Z\}\),
\[
 TQT^{-1}\not\sim_{\widetilde A}Q^{-1},
 \qquad
 TQT^{-1}\sim_{\widetilde A}\epsilon Q^{-1}.
 \tag{A.28.1}
\]
Each lift has order eight. More precisely the
unit exponents \(k\pmod8\) for which
\(TQT^{-1}\sim Q^k\) are exactly \(3,5\).
The same statements hold for either central
choice of each lift. All 1440 elements, the
central kernel, the 720 even conjugators and
the stated class tests are checked by the code;
the group order is not inferred solely from
the ATLAS label.\cite{AAtlasA6}

Let \(q_2:F\to\widetilde A\) send \(x,y\)
to \(X,Y\). If \(\Phi\) as in the statement
existed, the two homomorphisms
\(q_2\Phi\) and \(\operatorname{Ad}(T)q_2\)
would have the same projection to \(A_6\).
Their ratio is consequently a continuous central
character \(\eta:F\to\langle\epsilon\rangle\).
The multiplier-minus-one cusp conditions and
(A.28.1) require
\[
             \eta(x)=\eta(y)=\eta(z)=\epsilon.
\]
But \(xyz=1\), contradicting
\(\eta(x)\eta(y)\eta(z)=1\).
This proves the stronger assertion.
\end{proof}

The kernel of \(q\) is open normal; it is not
asserted to be characteristic. Any characteristic
refinement retaining this quotient and this action
inherits the displayed obstruction. Conversely,
we have not constructed an automorphism of such a
characteristic refinement or a relative pronilpotent
level inducing \(\vartheta\). The central calculation
explains why the finite inverse-class data alone
cannot be inserted into R2. The dyadic class
obstruction, rather than perfectness alone, is
what is retained here.

\statement{A.29}{Proved}{The new relative involution has exactly three section classes}
For the involution \(a_{12}\) constructed in A.25,
choose its actual representative \(\phi c\) on
\(F_{12}=F/R_{\mathrm{nil}}(U_{12})\).
Then
\[
                \bigl|H^1(C_2,F_{12})\bigr|=3,
\]
where the action is this specified representative.
Its relative dyadic component likewise has exactly
three section classes and has trivial fixed subgroup
in its orientation kernel. These are section
classes of the group extension, not conjugacy
classes of peripheral automorphisms.

\begin{proof}
First let
\(W=C_2\amalg C_2\amalg C_2=F\rtimes\langle c\rangle\),
and let \(U\) be any open normal subgroup of
\(W\) contained in \(F\). In
\(W/O^2(U)\), every nonidentity torsion element
of order two has nonidentity image in \(W/U\),
because \(U^{(2)}\) is torsion-free. If that
image generates \(C\cong C_2\), let
\(J\leq W\) be its inverse image. The open
Kurosh theorem expresses \(J\) as a free
profinite product of order-two factors and a
free profinite group. Since \(J/U=C_2\),
\[
                     J/O^2(U)=J^{(2)}.
\]
Indeed the left side is pro-2, and every finite
2-group quotient of \(J\) kills \(O^2(U)\).
The torsion theorem for the resulting free pro-2
product shows that the given involution is conjugate
to an image of a full factor of \(W\).
The three full factors remain nonconjugate,
because \(W^{\mathrm{ab}}=C_2^3\) survives this
quotient. Conjugacy under the even subgroup
is the same as full conjugacy for odd involutions:
one may multiply an odd conjugator by its
source involution to make it even.

The same open-product argument proves that every
such torsion centralizer is exactly the cyclic
subgroup it generates. Explicitly, if a centralizing
element has finite image \(g\), conjugating its
factor by a lift of \(g\) gives two factors of
\(J\) that become conjugate in \(J^{(2)}\).
The free pro-2 torsion theorem forces equality
of their Kurosh labels, so \(g\in C\).
The centralizer computation now takes place inside
\(J^{(2)}\), where the free-factor centralizer
is that factor. Its intersection with the even
subgroup is trivial. This uses the free-factor
case of Herfort--Ribes, not a closed-subgroup
Kurosh assertion.\cite{A_HR1985}

Apply this to \(U=U_{12}\), which is preserved
by the actual Galois action used in A.25.
The dyadic involution \(\Phi_\sigma c\) is
Galois-conjugate to \(c\), so its section set
has the same three classes and trivial even
centralizers. The choice of conjugating Galois
element preserves \(U_{12}\), by R3's quotient
covariance.

At each odd prime the kernel above the common
finite quotient is pro-\(p\). Every involution
in the specified quotient coset has an involution
lift, and any two such lifts with the same finite
image are conjugate by that kernel, by
Schur--Zassenhaus. Moreover the centralizer of
an involution surjects onto the centralizer of
its finite image: conjugate a lift by any chosen
finite centralizer lift, then apply the same
complement-conjugacy assertion in the kernel.

It follows that projection to the dyadic component
induces a bijection on section-conjugacy classes
for the complete fibre product. Surjectivity
follows by choosing odd-prime lifts with the
same finite image. For injectivity, choose a
dyadic conjugator and its finite image; at every
odd prime use centralizer surjectivity to adjust
a conjugator to this very same finite image.
The resulting tuple is an element of the fibre
product and conjugates the two sections. This
proves the asserted cardinality three.
\end{proof}

The unmarked pro-2 recognition theorem is fully
consistent with A.29. Herfort--Zalesskii's
Theorem~6.2 conjugates a finite 2-subgroup of
\(\operatorname{Aut}(F^{(2)})\) into the
automorphisms of a dense abstract free group.
It does not retain three prescribed peripheral
classes.\cite{AHerfortZalesskii2013}
Thus neither that theorem nor the section count
proves the peripheral conjugacy statement left
open in A.12. The actual involution of A.25
already passes the three-section test.

\statement{A.30}{Proved}{The level-twelve frames can also satisfy the involution relations}
The construction in A.25 can be made with its literal
relative frames \(A,\beta\in\Gamma\) so that the
associated r-normalizer and s-normalizer in
\(\Gamma\rtimes\langle\phi c\rangle\) are actual
subgroups \(S_3\) and \(C_2\times C_2\), with the
finite frame values still
\[
                         q(A)=1,\qquad q(\beta)=b.
\]
Consequently an R2 lift \(a=\widetilde\sigma c\)
can be chosen with
\[
 a^2\in L_{12}^{\mathrm{fr}},\qquad
 L_{12}^{\mathrm{fr}}
 =\{k\in\mathcal H:\Phi_k|_\Gamma=1,
                  \ A_k|_\Gamma=\beta_k|_\Gamma=1\}.
 \tag{A.30.1}
\]
This is a smaller kernel than the unframed
\(L_{12}\). It does not assert a lift at a
smaller characteristic level or solve the remaining
norm equation.

\begin{proof}
Write \(a_2=\Phi_\sigma c\) for the genuine
Galois involution used at the dyadic component,
and write \(A_2,\beta_2\) for its underlying
multiplier-one element's frames. Put
\[
 J_2=\Phi_\sigma(s)A_2,
 \qquad D_2=J_2^{-1}a_2,
 \qquad T_2=\beta_2^{-1}a_2.
\]
The element \(D_2\) inverts \(r\), because
\(c(r)=sr^{-1}s\). In the full Galois extension,
its square lies in \(C_{\widehat B}(r)=\langle r\rangle\)
and is fixed by inversion; hence \(D_2^2=1\).
Similarly \(T_2\) centralizes \(s\) and its
square lies in \(\langle s\rangle\). That square
cannot be \(s\): the abelianization of
\(\widehat B\rtimes\langle\Phi_\sigma c\rangle\)
is \(C_2^2\), and \(s\) has nonzero image.
Thus \(T_2^2=1\). These equations descend to
\(\Gamma_2\rtimes\langle a_2\rangle\).

In the common finite extended quotient, put
\(d=sc\) and \(z_0=[5I]\). Since
\(q\Phi_\sigma=q\), the frame values give
\[
                       q(D_2)=z_0d,
                       \qquad q(T_2)=b^{-1}c.
\]
Choose a pure \(h\in\Gamma_2\) with finite
image \([B]\), where \(B\) is the matrix of
A.25. Direct multiplication gives
\[
             d(B)=B^{-1},\qquad B^2=5I.
\]
It follows that the involution \(hD_2h^{-1}\)
has finite image \(d\). The inverse image of
\(\langle q(r),d\rangle\cong S_3\) in this
relative extension has pro-2 kernel. A.16 supplies
a section retaining this exact involution. Its
order-three subgroup is conjugate to \(\langle r\rangle\)
by an element of \(U_{12}^{(2)}\), by coprime
complement conjugacy. After this conjugation the
section has an involution \(D'_2\) satisfying
\[
        D'_2r(D'_2)^{-1}=r^{-1},\qquad q(D'_2)=d.
\]
Then \(k_2=D_2D'_2\) centralizes \(r\), belongs
to \(\Gamma_2\), and has finite image \(z_0\).
Replace \(A_2\) by \(A'_2=A_2k_2\).
The r-frame now has finite image one, and its
associated r-normalizing reflection is \(D'_2\).
At every odd prime retain the identity r-frame
and the standard r-normalizing reflection \(sc\).
Thus the r-frames glue, and their normalizers
glue to an actual finite \(S_3\).

At two retain \(\beta_2\) and \(T_2\).
For an odd prime \(p\), A.23 gives a surjection
\[
             C_{\Gamma_p}(s)\longrightarrow C_{Q_{12}}(q(s))
\]
with pro-\(p\) kernel, equivariant for \(c\).
The element \(b^{-1}c\) has order two in the
finite semidirect product, since \(c(b)=b^{-1}\).
Schur--Zassenhaus supplies an involution above it
in \(C_{\Gamma_p}(s)\rtimes\langle c\rangle\).
Write this involution as \(\beta_p^{-1}c\).
Then \(q(\beta_p)=b\), and \(\beta_p\) is a
frame for the identity action at that prime.
These frames glue with \(\beta_2\); their
involutions glue to \(T=\beta^{-1}(\phi c)\),
which centralizes \(s\). Since \(T\) is outside
\(\Gamma\), it and \(s\) generate \(C_2^2\).
The unchanged \(\tau\)-frame is one.

Apply R2 to these actual frames, and call the full
lift \(a=\widetilde\sigma c\). It still has
\(\Delta(a)=1\) and the same relative action as
A.25. For its s-frame, applying \(a\) twice
gives the pure conjugator \(a(\beta)\beta\).
Its relative image is one, exactly because
\((\beta^{-1}(\phi c))^2=1\).
For the r-frame put \(J=\Phi_{\widetilde\sigma}(s)A\).
Applying \(a\) twice gives the pure conjugator
\(a(J)J\), whose relative image is one because
\((J^{-1}(\phi c))^2=1\). Uniqueness of the full
pure frames identifies these with the frames of
\(a^2\). Together with its identity relative
action, this proves (A.30.1).
\end{proof}

The kernel \(L_{12}^{\mathrm{fr}}\) is invariant
under conjugation by this chosen \(a\). One may
verify this by composing actual peripheral actions
and their literal conjugators: an identity action
with both conjugators equal to one remains such
under framed conjugation. In the specified
\emph{framed} fibre, an order-two full lift exists
exactly when
\[
             v(ava^{-1})=a^{-2}
             \quad\text{for some }v\in L_{12}^{\mathrm{fr}}.
 \tag{A.30.2}
\]
Solving (A.30.2) would solve Lead~3. Failure to
solve it would exclude only this framed fibre;
a full involution over the unframed \(a_{12}\)
could still use different relative frames, and
is governed by A.25's equation in \(L_{12}\).
Neither equation is declared solvable or insoluble.
